\pdfbookmark[0]{English title page}{label:titlepage_en}
\aautitlepage{%
  \englishprojectinfo{
    Enhanced Person Re-Identification %title
  }{%
    Scientific Theme %theme
  }{%
    Fall Semester 2025 %project period
  }{%
    1 % project group
  }{%
    %list of group members
    Anne Sofie Ingerslev\\ 
    Lærke Stein Raaschou\\
    Jonas Holm Ingvorsen\\
    Stoyan Dimitrov Mihaylov\\
    Peter Gravgaard Andersen\\
    Rasmus Mølby Nielsen
    
  }{%
    %list of supervisors
    Andreas Aakerberg 
  }{%
    1 % number of printed copies
  }{%
    \today % date of completion
  }%
}{%department and address
  \textbf{The Technical Faculty of IT and Design}\\
  Aalborg University\\
  \href{http://www.aau.dk}{http://www.aau.dk}
}{% the abstract
  This project focuses on improving \acs{ReID} performance in edge case scenarios where poor resolution makes identification challenging, which is particularly relevant in law enforcement investigations. The main hypothesis is that enhancing low-resolution images through \acs{SR} before passing them to a \acs{ReID} network could improve identification accuracy. To investigate this, the project aims to design, implement, and deploy a proof-of-concept system that integrates \acs{SR} and \acs{ReID} technologies. The system is deployed as a containerized API service using Docker and FastAPI, enabling real-world testing and demonstrating the practical feasibility of the approach. This report presents the investigation, implementation, evaluation, and discussion of the integrated system, with the goal of understanding what can be done to improve identification accuracy on low-resolution surveillance footage and demonstrating how such a system can be deployed in practice.
 }

\cleardoublepage
